\section{Threats to Validity and Limitations}
\label{sec:threats}

\paragraph{Benchmark representativeness.}
The corpus comes from HeCBench and covers only the 95 benchmark programs that yielded the released 254-kernel evaluation subset.
That is a meaningful but still selective sample of CUDA and OpenMP GPU software.
The results should therefore be read as evidence about this benchmarked task, not as a claim about all GPU kernels or all accelerator software.

\paragraph{Hardware scope.}
We study only NVIDIA GPUs, and only four of them: RTX~3080, A10, A100, and H100.
The architectural generalization claims are therefore limited to those device families and their balance-point regimes.
The benchmark does not cover AMD or Intel GPUs, CPU-only kernels, or other accelerator programming models such as HIP and SYCL.

\paragraph{Task scope.}
The prediction target is deliberately narrow.
We predict only the first kernel invocation, only DRAM-level \rai, and only the per-precision FP16/FP32/FP64 decomposition used in this study.
We do not predict runtime, throughput, speedup, cache-aware arithmetic intensity, or end-to-end application performance.
Readers should therefore resist extrapolating the results beyond early first-invocation triage.

\paragraph{Ground-truth measurement.}
Ground truth comes from profiler counters, not from source-level reasoning.
That is the right choice for our task, but it means the denominator is DRAM-visible traffic rather than all source-visible loads and stores.
Cache write-back behavior, replay effects, and profiler noise can therefore influence the target quantity.
We reduce noise by averaging three first-invocation profiles, but this does not eliminate all measurement uncertainty.

\paragraph{Prompting and model stability.}
The released evaluation subset queries the expensive models once per completed sample.
We therefore cannot estimate run-to-run stochastic variability from the existing artifact, and closed-model behavior may drift across snapshots.
We address this by narrowing claims and by reporting completed-sample counts explicitly, but we cannot retroactively add repeated-query evidence.

\paragraph{Coverage imbalance.}
\gptoss does not complete the full evaluation subset: it has 858 source-only and 845 \sourcesass completed samples instead of 1{,}016.
The missing samples are primarily context-overflow or completion failures on long prompts, which is part of the deployment cost of cheaper long-context configurations.
Cross-model comparisons are therefore restricted to the 732 samples completed by all three models in both prompt settings.
This is a fairer comparison, but it also means that the cheapest model's reported quality is conditioned on the samples where it completed.

\paragraph{Baseline limitations.}
We now add context-only, deterministic, and lightweight learned tree baselines, and the learned baselines remain stable across five seeds.
But those baselines are still deliberately modest.
The deterministic baselines use approximate lexical and instruction-mix counts rather than compiler-grade analyzers.
The learned baselines are evaluated with grouped cross-validation on the released corpus rather than on a separately curated external training/test split.
Stronger compiler-driven or larger-scale learned baselines could therefore do better than the references reported here.

\paragraph{Failure-analysis validity.}
The feature/error analysis is exploratory because its source-feature labels come from an auxiliary LLM-voting pipeline, not human annotation.
We therefore treat the appendix heatmap and its prevalence/BH-corrected sanity check only as descriptive evidence about recurring error patterns.
It should not be interpreted as statistically validated causal evidence.

\paragraph{Possible benchmark leakage.}
Because the models are off-the-shelf and trained on large web corpora, some kernels or their close variants may have appeared in pretraining data.
We cannot rule out such contamination completely.
That is another reason to frame the result as an empirical benchmark study rather than a claim about intrinsic reasoning ability in isolation.
